\section{Discussion and Limitations}
\label{sec:discussion}

\subsection{Spatial Data Integration and Map Completeness}

The results demonstrate the importance of evaluating
geometric consistency and spatial coverage before using
historical soil cartography in contemporary GIS analysis.
Within the 178.47~km$^2$ Area of Interest (AOI), the
soil polygons covered approximately 162.852~km$^2$,
measured as their geometric union in EPSG:3795.
Of the 15.62~km$^2$ not covered by soil polygons,
approximately 11.42~km$^2$ corresponded to mapped
water surfaces, leaving 4.20~km$^2$ of non-water gaps.

Distinguishing water surfaces from unresolved gaps
prevents the absence of terrestrial soil polygons
over the reservoir from being incorrectly interpreted
as a soil-mapping defect. The remaining gaps require
examination against the original cartography and
other authoritative spatial sources before they
can be classified or filled.

Topology checks identified 90 polygon pairs with
positive-area intersections. Their cumulative
pairwise overlap area was approximately
0.054919~km$^2$. Although small relative to the
mapped extent, these overlaps can introduce
duplicate area contributions in spatial summaries
and overlay operations.

The original geometries were preserved to maintain
data provenance. Consequently, the processed dataset
should not be described as fully topologically
corrected. Area statistics based on sums of individual
polygons must be distinguished from statistics based
on their geometric union. Furthermore, cumulative
pairwise overlap areas do not necessarily equal the
geographic area affected by overlaps when three or
more polygons intersect.

Area calculations also varied slightly with the
coordinate reference system (CRS). In EPSG:3795,
the sum of individual polygon areas was
162.907~km$^2$, while their geometric union covered
162.852~km$^2$. After reprojection to EPSG:32617,
the corresponding values were 162.820~km$^2$
and 162.766~km$^2$, respectively. The union area
decreased by approximately 0.086~km$^2$
(0.053\%) after reprojection. In both CRSs,
the difference between the polygon-area sum and
the geometric union remained approximately
0.055~km$^2$, consistent with the contribution
of overlapping polygons to the summed area.
These results emphasize the need to specify the
CRS and area-calculation method when reporting
spatial coverage.

The integration of soil and vegetation polygons
provides a common spatial framework for examining
mapped attributes. Nevertheless, spatial coincidence
does not establish that a soil class determines
the observed vegetation condition. The source
layers may differ in production date, mapping
purpose, classification criteria, and positional
accuracy. These differences limit the interpretation
of direct spatial comparisons.

An additional area-accounting issue concerns the
vegetation dataset. The four reported vegetation
class areas sum to approximately 178.47~km$^2$,
whereas the stated vegetation coverage is
161.93~km$^2$. Their mutual exclusivity, possible
overlaps, and the corresponding area-calculation
procedures must therefore be verified before the
categories can be interpreted as an additive
partition of the AOI.


\subsection{Interpretation of Regional NDVI Dynamics}

The 15 retained Sentinel-2 observations show
substantial temporal variation within the selected
16.22~km$^2$ agricultural region.

Mean NDVI declined from approximately 0.688
on 3 November 2025 to 0.466 on 15 March 2026,
then varied during April and May before reaching
approximately 0.737 on 16 June 2026.

The boxplots, density curves, and spatial maps
indicate that these changes involved both
regional central tendency and the spatial
distribution of NDVI values.

These observations are consistent with changes
in vegetation conditions across an actively
managed agricultural landscape, but they do
not independently identify their causes.

The selected region contains multiple agricultural
polygons rather than an independently verified
single-crop planting area. Consequently, its
aggregate trajectory cannot be interpreted as
a potato growth curve.

Changes in crop composition, planting and
harvesting schedules, vegetation regrowth,
weeds, moisture conditions, and residual
image contamination could contribute to
the observed signal.

In particular, the June increase cannot be
attributed to a specific crop or management
practice without additional field evidence.

The regional analysis is descriptive rather
than causal. The association of the selected
polygons with mapped Ferrallitic soil units
provides geographic context, but the analysis
does not control for irrigation, agricultural
management, crop identity, or other potential
confounding factors.

The results therefore do not support claims
that one mapped soil class produced higher
NDVI or better crop performance than another.


\subsection{Field-Level Observations and Agricultural Context}

The three locally verified potato fields
provide a more specific spatial reference
than the broader agricultural region.

Their NDVI trajectories were not uniform.
Among the retained observations, the highest
mean NDVI occurred on 15 March 2026 in Field~1,
on 16 June 2026 in Field~2, and on
24 January 2026 in Field~3.

These differences illustrate why an aggregate
regional NDVI curve should not be interpreted
as representative of every individual field.

However, the dates of the highest observed
mean NDVI values do not necessarily represent
the actual seasonal vegetation maxima.
Unobserved peaks may have occurred between
the available satellite acquisitions.

Local verification supports the identification
of the sampling locations as potato fields,
but it does not establish that potatoes
occupied each sampled polygon continuously
throughout the November 2025--June 2026
observation period.

In particular, post-harvest NDVI may reflect
subsequent crops, vegetation regrowth, weeds,
or other vegetation rather than the original
potato crop.

The circular sampling polygons are
approximations rather than surveyed field
boundaries. Their geometry and placement
may therefore influence the extracted NDVI
statistics, particularly where sampling
areas intersect heterogeneous vegetation
or field boundaries.

Field-specific planting dates, harvesting
records, accurate field boundaries, and
repeated ground observations would be
needed to assign crop-development stages
and interpret field-level differences
with greater confidence.


\subsection{Remote-Sensing and Reproducibility Limitations}

The regional quality filter retained
acquisitions with less than 10\% cloud
coverage according to the selected SCL
cloud classes.

This acquisition-level threshold reduces
potential contamination but does not
guarantee that every pixel used for NDVI
statistics is free from clouds, cirrus,
haze, or shadows.

Cloud shadows were evaluated separately,
and numerical NDVI validity alone does
not constitute a cloud mask.

Future analyses should implement and
document explicit pixel-level quality
masks, including cloud shadows and
other unsuitable SCL classes, and report
the number and proportion of usable
pixels for each acquisition.

The temporal sample is also irregular.
Fifteen retained dates span November
2025 through June 2026, with unequal
intervals between observations.

These data support comparisons among
the retained acquisition dates but do
not provide a continuously observed
vegetation-development curve.

Short-lived changes and seasonal
maxima may have occurred between
acquisitions. Furthermore, the nearly
constant number of numerically valid
pixels does not independently establish
equivalent atmospheric quality across
the observations.

Radiometric preprocessing used a
quantification factor of 10,000 and
an additive offset of $-1000$ for
both spectral bands, based on inspection
of 111 available product metadata files.

These parameters were also applied
to some images without individually
available XML metadata.

Although this assumption was documented,
product-specific verification remains
necessary to establish the consistency
of radiometric preprocessing across
the full time series.

Unlike a common multiplicative factor,
an additive radiometric offset does
not cancel during NDVI calculation.
Consequently, incorrect offset
parameters may affect the resulting
vegetation-index values.

Additional sources of uncertainty
include the 10~m spatial resolution
of the spectral bands, the 20~m
resolution of the SCL, resampling
choices, and potential positional
differences between historical
vector boundaries and satellite
imagery.

Overall, the principal contribution
of GeoSoil-Cuba is a documented and
reproducible workflow for integrating
historical soil cartography, existing
land-cover information, and
satellite-derived vegetation indices.

The results support spatial
data-quality assessment and
exploratory vegetation monitoring,
but they do not independently
validate soil properties, crop
yield, or causal relationships
between soil classes and
agricultural performance.

The principal priorities for
further development are targeted
topology review, harmonized area
accounting, product-specific
radiometric verification,
explicit pixel-level quality
masking, and field observations
aligned with satellite
acquisition dates.